\chapter{Boy Scouts Forever!}
I was a Boy Scout once, although I couldn't tell a slip knot form a granny knot, or a reef knot from a thief knot, expect that a thief, should you happen to catch one. I have never caught a thief and wouldn't know what to do with one since I can't tie a knot. Just let him go with a warning, I suppose. Tell him to become a Boy Scout.
\par
'Be prepared!' That's the Boy Scout motto. And a good one, too. But I never seen to be well prepared for anything, be in an exam or a journey or the roof blowing off my room. I get halfway through a speech and then forget what I have to say next. Or I make a new suit to attend a friend's and then turn up in my pyjamas.
\par
So how did I, the most impractical of boys, become a Boy Scout? I was at boarding school in Simla when it happened.
\par
Well, it seems a rumour had gone around the junior school (I was still a junior then) that I was a good cook. I had never cooked anything in my life, but I had spent a lot of time in the tuck shop making suggestions and advising samosas, jalebis, tikkis, and parkoras. For my unwanted advice, he would favour me with an occasional free samosa, so naturally I looked upon him as friend and benefactor. With this qualification, I was given a cookery badge and put in charge of our troop's supply of rations.
\par
There were about twenty of us in our troop, and during the summer break, our Scoutmaster, Mr Oliver, took us on a camping expedition to Tara Devi, a temple-crowned mountain a few miles outside Simla. That first night we were put to work, peeling potatoes, skinning onions, shelling peas and pounding masalas. These various ingredients being ready, I was asked-as the troop's cookery expert-what should be done with them. 'Put everything in the big \textit{degchi},' I ordered. 'Pour half a tin of ghee over the lot. Add some nettle leaves and cook for half and hour.'
\par
When this was done, everyone had a taste, but the general opinion was that the dish lacked something.
\par
'More salt,' I suggested.
\par
More salt was added. It still still lacked something. 'Add cup of a sugar,' I ordered.
\par
Sugar was added to the concoction. But still it lacked something.
\par
'We forgot to add tomatoes,' said Bimal, one of the Scouts.
\par
'Never mind,' I said. 'We have tomato sauce. And a bottle of tomato sauce!'
\par
'How about some vinegar?' asked another boy.
\par
'Just the thing,' I agreed.
\par
A cup of vinegar!'
\par
'Now it's too sour,' said one of the tasters.
\par
'What jam did we bring?' I asked.
\par
'Gooseberry jam.'
\par
'Just the thing. Empty the bottle!'
\par
The dish was a great success. Everyone enjoyed it, including Mr Oliver, who had no idea what went into it.
\par
'What's this called? he asked.
\par
'It's an all-Indian sweet-and-sour jam-potato curry,' I ventured.
\par
'For short, just call if a Bond-bhujji,' said Bimal. I had earned my cookery badge!
\vspace{0.5in}


Poor Mr Oliver! He wasn't really cut out to be a Scoutmaster, any more than I was meant to be a Scout. The following day, he announced that he would give us a lesson in tracking. He would take a ha;f-hour start and walk into the forest, leaving behind him a trail of broken twigs, chicken feathers, pine cones and chestnuts, and we were to follow the trail until we found him.
\par
Unfortunately, we were not very good trackers. We did follow Mr Oliver's trail some way into the forest, but were distracted by a pool of clear water which looked very inviting. Abandoning our uniforms, we jumped into the pool and had a great time romping around or just lying on the grassy banks and enjoying the sunshine. A couple of hours later, feeling hungry, we returned to our campsite and set about preparing the evening meal. Bond-bhujji again, but with further variations.
\par
It was growing dark, and we were beginning to worry about Mr Oliver's whereabouts when he limped into camp, assisted by a couple of local villagers. Having waited for us at the far end of the forest for a couple of hours, he had decided to return by following his own trail, but in the gathering gloom he was soon lost. Some locals returning from the temple took charge of him and escorted him back to camp. He was very angry and made us return all our good-conduct and other badges, which he stuffed into his haversack. I had to give up my cookery badge, too.
\par
An hour later, when we were all preparing to get into our sleeping bags for the night, Mr Oliver called out, 'Where's dinner?'
\par
'We've had ours,' said Bimal. 'Everything is finished, sir.'
\par
'Where's Bond? He's supposed to be the cook. Bond, get up and make me an omelette.'
\par
'Can't, sir.'
\par
'Why not?'
\par
'You have my badge. Not allowed to cook without it. Scout rule, sir.'
\par
'Never heard such an omelette before,' confessed Mr Oliver, blowing out his cheeks.
\par
'A little too hot, but otherwise quite interesting.'
\par
'Would you like another, sir?'
\par
'Tomorrow, Bond, tomorrow. We'll breakfast early tomorrow.'
\par
But we had to break up our camp very early the next day. In the early hours, a bear had strayed into our camp, entered the tent where our stores were kept and wreaked havoc with all our provisions, even rolling our biggest degchi down the hillside.
\par
In the confusion and uproar that followed, the bear entered Mr Oliver's tent (he was already outside, fortunately) and came out entangled in Mr Oliver's dressing gown. It then made off in the direction of the forest.
\par
A bear in a dressing gown? It was a comical sight. And through we were a troop of brave little Scouts, we thought it better to let the bear keep the gown.
